\documentclass[12pt]{article}
\usepackage[margin=0.6in,top=1.15in,bottom=0.55in]{geometry}
\usepackage{lmodern}
\usepackage{microtype}
\usepackage{titlesec}
\usepackage{enumitem}
\usepackage[hidelinks]{hyperref}
\usepackage{../../latex/grader-worksheet}
\setlength{\parindent}{0pt}
\setlength{\parskip}{0.25em}
\titleformat{\section}{\large\bfseries\scshape}{}{0pt}{}
\GraderSetup{assignment-id=U1L19LE,template-revision=1}
\begin{document}
\begin{center}
{\LARGE\bfseries Week 19 Lit Example Worksheet}\par\vspace{0.05em}
{\large\scshape Tragic Judgment and Disciplined Thought}\par\vspace{0.12em}
\rule{0.78\textwidth}{0.5pt}
\end{center}
\textbf{Purpose:} Build a miniature portfolio defense from Shakespeare: claim, textual evidence, warrant, qualification, and concrete transfer.

\section*{Part 1: Macbeth --- Evidence and Revision}
\textbf{Question 1.1}\\[-0.15em]
What observation challenges Macbeth's confidence, and how does his first response differ from his later response?
\GraderAnswerBox{Part 1 Q1}{2.2in}

\textbf{Question 1.2}\\[-0.15em]
Make a limited claim about disciplined judgment that this sequence supports. Name the Whately habit and state the warrant linking passage to claim.
\GraderAnswerBox{Part 1 Q2}{2.45in}

\newpage
\section*{Part 2: Macbeth --- Qualification and Transfer}
\textbf{Question 2.1}\\[-0.15em]
What does the passage not prove about logic or character? Add one qualification that keeps your claim from overreaching.
\GraderAnswerBox{Part 2 Q1}{2.35in}

\textbf{Question 2.2}\\[-0.15em]
Transfer the warning to a concrete nonliterary setting. Name the pressure and the action a disciplined thinker should take.
\GraderAnswerBox{Part 2 Q2}{2.45in}

\newpage
\section*{Part 3: Brutus --- Independent Defense}
\textbf{Question 3.1}\\[-0.15em]
Make a precise claim about the pressure point in Brutus's reasoning. Use words from the passage as evidence; do not summarize the plot.
\GraderAnswerBox{Part 3 Q1}{2.4in}

\textbf{Question 3.2}\\[-0.15em]
State the warrant connecting your evidence to one disciplined habit, then name a fair counterpressure or alternative the conclusion must answer.
\GraderAnswerBox{Part 3 Q2}{2.55in}

\newpage
\section*{Part 4: Comparative Closing}
\textbf{Question 4.1}\\[-0.15em]
Compare the two failures of judgment. What shared habit is missing, and how do the passages show different forms of pressure on that habit?
\GraderAnswerBox{Part 4 Q1}{2.55in}

\textbf{Question 4.2}\\[-0.15em]
Write a closing paragraph for a portfolio defense using one passage. Include a claim, evidence, warrant, qualification, and transfer beyond the logic unit.
\GraderAnswerBox{Part 4 Q2}{2.75in}
\end{document}
